\section{Tinjauan Kasus}
\label{sec:tinjauankasus}

Proyek yang dikerjakan selama periode magang di Badan Pendapatan Daerah (Bapenda) Kota Surabaya adalah "Pengembangan Sistem Monitoring Pendapatan Daerah (MonPD Reborn) Berbasis .NET". Proyek ini berfokus pada perancangan dan implementasi beberapa modul vital dalam sebuah sistem informasi internal yang bertujuan untuk meningkatkan efektivitas pengawasan, penagihan, dan analisis data perpajakan daerah.

Sistem ini dirancang untuk dapat diakses oleh pegawai internal Bapenda dengan tujuan utama menyediakan dashboard analitik dan tabel data yang interaktif. Website ini dibangun untuk mengolah dan memvisualisasikan data dari berbagai jenis pajak, seperti Pajak Reklame, Hotel, Restoran, dan lainnya. Tujuannya adalah untuk memberikan pandangan yang komprehensif dan real-time mengenai aktivitas wajib pajak (WP), realisasi target pendapatan, serta progres aktivitas lapangan seperti pemeriksaan dan penagihan, sehingga dapat mendukung pengambilan keputusan yang lebih cepat dan berbasis data.
